% colophon.tex
% ============================================================
%  صفحه‌ی پایانی (شناسنامه‌ی فنی)
% ============================================================

\clearpage
\thispagestyle{empty}

\begin{tikzpicture}[remember picture, overlay]
  \fill[PrimaryDark] (current page.south west) rectangle (current page.north east);
  
  % Content
  \node[text width=12cm, align=center, text=white]
    at (current page.center) {
      {\color{AccentGold}\rule{8cm}{0.5pt}}\\[20pt]
      {\Large\bfseries\color{AccentGold} شناسنامه‌ی فنّی کتاب}\\[20pt]
      {\color{AccentGold}\rule{8cm}{0.5pt}}\\[30pt]
      {\small
        \begin{tabular}{rl}
          \textcolor{AccentGold}{\bfseries عنوان:} &
            دولت، ملت و تنوع قومی \\[6pt]
          \textcolor{AccentGold}{\bfseries زیرعنوان:} &
            مبانی نظری، الگوهای تطبیقی و راهنمای اجرایی برای ایران \\[6pt]
          \textcolor{AccentGold}{\bfseries تألیف:} &
            پژوهشگر سیاست تطبیقی \\[6pt]
          \textcolor{AccentGold}{\bfseries حروف‌چینی:} &
            \en{\XeLaTeX{} + TikZ + tcolorbox} \\[6pt]
          \textcolor{AccentGold}{\bfseries قلم فارسی:} &
            \en{XB Niloofar / IRanSans} \\[6pt]
          \textcolor{AccentGold}{\bfseries قلم لاتین:} &
            \en{Libertinus Serif / Sans} \\[6pt]
          \textcolor{AccentGold}{\bfseries ویرایش:} &
            اول — بهار ۱۴۰۴ \\[6pt]
          \textcolor{AccentGold}{\bfseries تعداد فصول:} &
            ۶ فصل + پیوست‌ها \\[6pt]
          \textcolor{AccentGold}{\bfseries قطع:} &
            وزیری (\en{A4}) \\[6pt]
          \textcolor{AccentGold}{\bfseries شابک:} &
            \en{978-600-XXX-XXX-X}
        \end{tabular}
      }\\[30pt]
      {\color{AccentGold}\rule{8cm}{0.5pt}}\\[15pt]
      {\footnotesize\color{white!70}
        این کتاب با نرم‌افزارهای آزاد حروف‌چینی شده\\
        و تمامی نمودارها و تصاویر با \en{TikZ/PGF} تولید شده‌اند.\\[8pt]
        بازنشر با ذکر منبع بلامانع است.
      }
    };
\end{tikzpicture}

\clearpage